\subsection{Quadratic Probing}
\begin{frame}{Quadratic Probing}
\[
h(k,i)=\bigl(h'(k)+c_1i+c_2i^2\bigr)\bmod m.
\]
원본 단순형은 $m=13$, $c_1=0$, $c_2=1$:
\[
h(k,i)=\bigl(h'(k)+i^2\bigr)\bmod13.
\]
\httakeaway{offset가 $0,1,4,9,16,\ldots$로 커져 contiguous primary run의 영향을 완화한다.}
\end{frame}
\begin{frame}{Quadratic Trace: 30 삽입}
\centering
\begin{tabular}{c|cccccc}\toprule
입력&15&18&43&37&45&30\\
home&2&5&4&11&6&4\\
final&2&5&4&11&6&8\\\bottomrule
\end{tabular}
\vspace{5mm}
\begin{tikzpicture}[x=.8cm]
\foreach \i in {0,...,12}{\node[ht empty,minimum width=8mm](s\i)at(\i,0){E};\node[font=\tiny,text=AlgoGray]at(\i,-.5){\i};}
\foreach \i/\v in {2/15,4/43,5/18,6/45,11/37}{\node[ht occupied,minimum width=8mm]at(s\i){\v};}
\only<1|handout:0>{\node[ht collision,minimum width=8mm]at(s4){43};\node[callout]at(6,1.2){$i=0$: $4+0=4$};}
\only<2|handout:0>{\node[ht collision,minimum width=8mm]at(s5){18};\node[callout]at(6,1.2){$i=1$: $4+1=5$};}
\only<3|handout:1>{\node[ht probe,minimum width=8mm]at(s8){30};\node[callout]at(6,1.2){$i=2$: $4+4=8$};}
\end{tikzpicture}
\end{frame}
\begin{frame}{Coverage는 자동 보장이 아니다}
\[
(h'+i^2)\bmod13,\quad i=0,1,\ldots
\]
\centering\begin{tabular}{c|rrrrrrrr}\toprule
$i$&0&1&2&3&4&5&6&7\\
$i^2\bmod13$&0&1&4&9&3&12&10&10\\\bottomrule
\end{tabular}
\begin{alertblock}{중요}$i=6$과 $i=7$부터 repeated offset이 보인다. $m,c_1,c_2$와 load에 따라 일부 slot을 방문하지 않아 table이 덜 찼어도 insertion이 실패할 수 있다.\end{alertblock}
\end{frame}
\begin{frame}{Secondary Clustering}
\centering\begin{tikzpicture}[node distance=7mm]
\node[hash box](a){15: home 2};
\node[hash box,below=of a](b){28: home 2};
\node[hash box,below=of b](c){41: home 2};
\node[hash box,right=25mm of b](p){shared sequence\\$2,3,6,11,5,\ldots$};
\draw[probe arrow](a)--(p);\draw[probe arrow](b)--(p);\draw[probe arrow](c)--(p);
\end{tikzpicture}
\begin{block}{정의}같은 initial hash value를 갖는 key들이 동일한 quadratic probe sequence를 따라가는 현상이다.\end{block}
\end{frame}
\begin{frame}{Primary vs Secondary Clustering}
\small\centering\begin{tabular}{lll}\toprule
&Primary&Secondary\\\midrule
대표 방식&linear probing&quadratic probing\\
합쳐지는 key&서로 다른 home도 가능&같은 home\\
원인&contiguous occupied run&동일 offset sequence\\
결과&run의 자기 강화&같은 home끼리 같은 경로\\\bottomrule
\end{tabular}
\end{frame}
\begin{frame}{Quadratic Probing Checkpoint}
\begin{enumerate}\item 30의 probe indices는?\item quadratic probing이 모든 slot을 방문한다고 말할 수 있는가?\item secondary clustering의 정확한 정의는?\end{enumerate}
\begin{exampleblock}{답}$4,5,8$; parameter/capacity policy 없이는 아니다; 같은 home key들이 같은 quadratic sequence를 공유하는 현상.\end{exampleblock}
\end{frame}
